\section{Билет 63. Тепловая машина. КПД тепловой машины. Вычисление КПД цикла Дизеля. Степень сжатия и расширения.}

\begin{definition}
\textbf{Тепловой машиной} называется устройство, преобразующее внутреннюю энергию топлива в механическую работу с помощью кругового процесса. \textbf{КПД} машины определяется как отношение полезной работы за цикл к полученной от нагревателя теплоте:
\begin{equation}
    \eta = \frac{A}{Q_1} = 1 - \frac{|Q_2|}{Q_1},
\end{equation}
где $Q_1$ --- теплота, полученная от нагревателя, $|Q_2|$ --- теплота, отданная холодильнику.
\end{definition}

\begin{definition}
\textbf{Цикл Дизеля} --- идеализированный термодинамический цикл дизельного двигателя внутреннего сгорания. Состоит из четырёх процессов:
\begin{enumerate}
    \item 1--2: Адиабатическое сжатие воздуха (температура повышается до $T_2$);
    \item 2--3: Изобарное подведение теплоты $Q_1$ (впрыск и сгорание топлива);
    \item 3--4: Адиабатическое расширение (рабочий ход);
    \item 4--1: Изохорное отведение теплоты $|Q_2|$ (выпуск отработавших газов).
\end{enumerate}
\end{definition}

\begin{theorem}[Вывод КПД цикла Дизеля]
Для изобарного и изохорного процессов количество теплоты:
\begin{equation}
    Q_1 = \nu C_p (T_3 - T_2), \quad |Q_2| = \nu C_V (T_4 - T_1).
\end{equation}
КПД выражается через температуры:
\begin{equation}
    \eta_{\text{Д}} = 1 - \frac{1}{\gamma} \frac{T_4 - T_1}{T_3 - T_2}, \label{eq:diesel_eff_temp}
\end{equation}
где $\gamma = C_p/C_V$ --- показатель адиабаты.
\par\noindent Введём геометрические параметры цилиндра:
\begin{itemize}
    \item \textbf{Степень сжатия} $n = V_1/V_2$;
    \item \textbf{Степень предварительного расширения} (степень отсечки) $m = V_3/V_2$.
\end{itemize}
Используя уравнения процессов:
\begin{itemize}
    \item Адиабата 1--2: $T_2 = T_1 n^{\gamma-1}$;
    \item Изобара 2--3: $\frac{T_3}{T_2} = \frac{V_3}{V_2} = m \Rightarrow T_3 = T_1 m n^{\gamma-1}$;
    \item Адиабата 3--4: $T_4 = T_3 \left(\frac{V_3}{V_4}\right)^{\gamma-1}$. Так как $V_4 = V_1$, то $\frac{V_3}{V_4} = \frac{m V_2}{n V_2} = \frac{m}{n}$. Следовательно, $T_4 = T_1 m n^{\gamma-1} \left(\frac{m}{n}\right)^{\gamma-1} = T_1 m^\gamma$.
\end{itemize}
Подставляя $T_2, T_3, T_4$ в \eqref{eq:diesel_eff_temp} и сокращая на $T_1$, получаем:
\begin{equation}
    \eta_{\text{Д}} = 1 - \frac{1}{n^{\gamma-1}} \frac{m^\gamma - 1}{\gamma(m - 1)}. \label{eq:diesel_eff_final}
\end{equation}
\end{theorem}

\begin{property}
\textbf{Анализ формулы \eqref{eq:diesel_eff_final}:}
\begin{enumerate}
    \item КПД растёт с увеличением степени сжатия $n$ и уменьшением степени отсечки $m$.
    \item При $m \to 1$ (изобарное подведение теплоты отсутствует, цикл вырождается в цикл Отто) формула переходит в $\eta_{\text{Отто}} = 1 - 1/n^{\gamma-1}$.
    \item Реальные дизельные двигатели имеют $n \approx 14 \text{--} 22$, $\gamma \approx 1.4$, что даёт теоретический $\eta_{\text{Д}} \approx 0.60 \text{--} 0.68$. Реальный КПД составляет $35 \text{--} 45\%$.
    \item В отличие от бензиновых двигателей, дизели сжимают только воздух, что позволяет использовать высокие степени сжатия без риска детонации, обеспечивая более высокий КПД.
\end{enumerate}
\end{property}
